\documentclass[a4paper,10.5pt]{article}
\usepackage[utf8]{inputenc}
\usepackage[T1]{fontenc}
\usepackage[french]{babel}
\usepackage[left=1cm,right=1cm,top=.5cm,bottom=0cm]{geometry}
\usepackage{enumitem}
\usepackage{hyperref}
\usepackage{xcolor}
\usepackage{titlesec}
\usepackage{fontawesome5}
\usepackage{lmodern}
\usepackage[normalem]{ulem}

% --- Couleurs EDF ---
\definecolor{primary}{RGB}{0, 85, 164}
\definecolor{secondary}{RGB}{80, 80, 80}

% --- Config Liens ---
\hypersetup{colorlinks=true, linkcolor=primary, filecolor=primary, urlcolor=primary}

% --- Titres ---
\titleformat{\section}{\large\bfseries\color{primary}\uppercase}{}{0em}{}[\titlerule]
\titlespacing{\section}{0pt}{8pt}{5pt}

% --- Commandes ---
\newcommand{\resumeItem}[1]{\item\small{#1}}

\begin{document}

% --- EN-TÊTE ---
\begin{center}
    {\Large \textbf{Loïc Aron MBASSI EWOLO}} \\ \vspace{4pt}
    \textbf{\large Alternance Data Scientist IA} \\
    \textit{\small Disponible dès septembre 2026 pour 12 mois} \\ \vspace{4pt}
    \small
    \faMapMarker\ Cergy, France (Mobile IDF) \quad
    \faPhone\ (+33) 7 59 52 33 98 \quad
    {\color{primary}\faEnvelope\ \href{mailto:loicaron.mbassiewolo@ensea.fr}{loicaron.mbassiewolo@ensea.fr}} \\ \vspace{2pt}
    {\color{primary}\faLinkedin\ \href{https://www.linkedin.com/in/loic-mbassi/}{linkedin.com/in/loic-mbassi}} \quad
    {\color{primary}\faGithub\ \href{https://github.com/Nameless0l}{github.com/Nameless0l}} \quad
    {\color{primary}\faGlobe\ \href{https://mbassiloic.tech}{mbassiloic.tech}}
\end{center}

\vspace{-6pt}

% --- PROFIL ---
\noindent\small{Conception et déploiement d'un système multi-agents (RAG, LangChain, Google ADK) orchestrant des LLM pour des recommandations contextuelles via APIs. Formation double diplôme ENSEA/Polytechnique Yaoundé en IA et mathématiques appliquées. Expérience en recherche IA (MILA, Québec) et en NLP (Transformers, BERT, GPT-2). Compétences Python solides, appétence pour l'IA agentique et les architectures LLM.}

% --- FORMATION ---
\section{Formation}
\begin{itemize}[leftmargin=*, label=\tiny$\bullet$, nosep]
    \item \textbf{ENSEA -- École Nationale Supérieure de l'Électronique et de ses Applications} \hfill \textit{2025 -- 2027} \\
    \small{Double diplôme d'Ingénieur -- Deep Learning, Traitement du Signal, IA, Systèmes Embarqués.}
    \item \textbf{École Polytechnique de Yaoundé (1\textsuperscript{ère} école d'ingénieur CEMAC)} \hfill \textit{2021 -- 2025} \\
    \small{Diplôme d'Ingénieur de Conception (Master 2) : \textbf{Mathématiques Appliquées}, Algorithmique, Génie Logiciel.}
\end{itemize}

% --- COMPÉTENCES ---
\section{Compétences Techniques}
\begin{itemize}[leftmargin=*, nosep]
    \item \small{\textbf{IA \& LLM :} Python, PyTorch, TensorFlow, LangChain, RAG, Google ADK, Transformers (BERT, GPT-2), Prompt Engineering.}
    \item \small{\textbf{Data Science :} Pandas, NumPy, Scikit-learn, Matplotlib, Seaborn, Streamlit, analyse de séries temporelles.}
    \item \small{\textbf{Architecture \& Outils :} FastAPI, Flask, Docker, Git, API REST, RabbitMQ, MongoDB, MySQL.}
    \item \small{\textbf{Mathématiques :} Algèbre linéaire (SVD), Probabilités/Statistiques, Calcul différentiel, Analyse numérique.}
    \item \small{\textbf{Langues :} Français (natif), Anglais (professionnel/technique).}
\end{itemize}

% --- EXPÉRIENCES / PROJETS ---
\section{Expériences \& Projets IA}
\begin{itemize}[leftmargin=*, label={-}]

\item \textbf{Système Multi-Agents Agricole (Recherche)} \hfill \textit{2024 -- 2025} \\
\textit{Projet académique -- IA agentique} \hfill \textit{Google ADK, RAG, LangChain, Gemini, FastAPI, Docker}
\vspace{-2pt}
    \begin{itemize}[leftmargin=0.4cm, nosep, label=\tiny$\bullet$]
        \item \small{Conception d'un système de 4 agents IA collaboratifs utilisant Gemini comme moteur LLM et RAG pour l'enrichissement contextuel.}
        \item \small{Orchestration via Google ADK : raisonnement, planification, résolution de conflits entre agents spécialisés.}
        \item \small{Intégration d'APIs externes (OpenWeather, marchés) pour recommandations agricoles multi-critères.}
        \item \small{Déploiement Python (Poetry, FastAPI), containerisation Docker.}
    \end{itemize}
\vspace{3pt}

\item \textbf{Assistant de Recherche @ MILA (Institut Québécois d'IA)} \hfill \textit{Juin 2025 -- Sept 2025} \\
\textit{Remote -- Montréal, Canada} \hfill \textit{Python, PyTorch, TensorFlow}
\vspace{-2pt}
    \begin{itemize}[leftmargin=0.4cm, nosep, label=\tiny$\bullet$]
        \item \small{Étude du phénomène ``Grokking'' (généralisation tardive des réseaux de neurones) sur des MLP.}
        \item \small{Reproduction d'expériences état de l'art, analyse mathématique de la convergence.}
        \item \small{Pipelines de tests statistiques, visualisation et documentation rigoureuse des résultats.}
    \end{itemize}
\vspace{3pt}

\item \textbf{Stage Recherche -- Labo ICS, Florence (Italie)} \hfill \textit{Avr 2026 -- Août 2026} \\
\textit{Analyse de Signaux \& Neurosciences Computationnelles} \hfill \textit{Python, MATLAB, Traitement du signal}
\vspace{-2pt}
    \begin{itemize}[leftmargin=0.4cm, nosep, label=\tiny$\bullet$]
        \item \small{Analyse de signaux et développement d'outils en Python pour la recherche en neurosciences computationnelles.}
        \item \small{Lecture et exploitation de publications scientifiques en anglais, proposition d'axes de recherche.}
        \item \small{Maintenance et évolution de bibliothèques MATLAB pour le laboratoire.}
    \end{itemize}
\vspace{3pt}

\item \textbf{1\textsuperscript{er} Prix IndabaX Africa 2025 -- IA Santé} \hfill \textit{2025} \\
\textit{Hackathon Pan-Africain} \hfill \textit{Python, Pandas, Flask, Streamlit}
\vspace{-2pt}
    \begin{itemize}[leftmargin=0.4cm, nosep, label=\tiny$\bullet$]
        \item \small{Nettoyage et structuration de données médicales, déploiement d'API Flask.}
        \item \small{Dashboard Streamlit interactif pour visualisation des KPIs santé.}
    \end{itemize}
\vspace{3pt}

\item \textbf{Détection de contenus sexistes -- Women Safe (NLP)} \hfill \textit{2023} \\
\textit{Compétition MLPC} \hfill \textit{BERT, GPT-2, Transformers, Python}
\vspace{-2pt}
    \begin{itemize}[leftmargin=0.4cm, nosep, label=\tiny$\bullet$]
        \item \small{Modèle NLP de détection de misogynie dans les textes (réseaux sociaux, offres d'emploi).}
        \item \small{Fine-tuning de modèles Transformers (BERT, GPT-2), analyse des biais éthiques.}
    \end{itemize}
\vspace{3pt}

\end{itemize}

% --- DISTINCTIONS ---
\section{Distinctions \& Leadership}
\begin{itemize}[leftmargin=*, label=\tiny$\bullet$, nosep]
    \item \small{\textbf{1\textsuperscript{er} Prix} CITS National Hackathon 2024 (75 équipes) -- Optimisation collecte déchets (ML, algorithme du voyageur de commerce).}
    \item \small{\textbf{Lead AI Cell} (ENSPY) : +50\% membres, collaboration avec mentors Google DeepMind.}
    \item \small{\textbf{Workshops animés :} ML, Fine-tuning LLM, React, Laravel. Articles techniques sur Medium.}
    \item \small{\textbf{Certif.} Supervised ML (Stanford/DeepLearning.AI), Python for Data Science (IBM/Coursera).}
\end{itemize}

\end{document}
